% !TeX program = xelatex
\documentclass[11pt]{article}

% Self-contained, publisher-neutral Chinese cover letter.
\usepackage[margin=1in]{geometry}
\usepackage{fontspec}
\usepackage{xeCJK}
\usepackage{microtype}
\usepackage{xcolor}
\usepackage[hidelinks]{hyperref}
\input{author_metadata.tex}

\definecolor{ink}{HTML}{20262E}
\definecolor{placeholder}{HTML}{6B7280}

\color{ink}
\setlength{\parindent}{0pt}
\setlength{\parskip}{0.72\baselineskip}
\pagestyle{empty}

% 灰色文字为写作提示；正式投稿前必须替换或删除。
\newcommand{\guidance}[1]{\textcolor{placeholder}{【#1】}}

\begin{document}

\hfill\today\par
\EditorName\par
\EditorTitle\par
\textit{\JournalName}

尊敬的编辑：

现随信提交题为《\ManuscriptTitle》的稿件，拟以\ArticleType 的形式投稿至
\textit{\JournalName}，恳请贵刊审议。

%%COVER_BODY%%

感谢您审阅本稿件。期待您的回复。

此致\par
敬礼！\par
\vspace{0.22in}

\begingroup
\setlength{\parskip}{2pt}
{\bfseries\CorrespondenceAuthorNameZh}\par
通讯作者，谨代表全体作者\par
\CorrespondenceAuthorAffiliationZh\par
\CorrespondenceAuthorAddress\par
电子邮箱：\href{mailto:\CorrespondenceAuthorEmail}{\CorrespondenceAuthorEmail}
\endgroup

\end{document}
